% Generated by scripts/blueprint_gen.py; do not edit.
\chapter{The Dolbeault sheaf and the main theorem}\label{chap:F}
\textbf{Setting.} E a finite-dimensional complex normed space, n = dim\textsubscript{ℂ} E, \texttt{E : Type}. F a finite-dimensional complex normed space of values (F = ℂ for the main theorem; general F is for Z3 and Z8). M a complex manifold modelled on E (\texttt{ChartedSpace E M}, \texttt{IsManifold 𝓘(ℂ, E) ∞ M}), not necessarily compact. Only F = ℂ is needed on the critical path; F2–F4 and F7 are stated for general F because it is free there. Conventions §6 fixes all normalizations.

\begin{theorem}[M — Main theorem conditional]
\label{M}
\lean{Cohesive.MainStatementConditional, main_theorem_conditional}
\leanok
\uses{C1, C2, F10, F12, H1, H2, S9}
Let E be a finite-dimensional complex normed space and X a compact Hausdorff complex manifold modelled on E. Let A = 𝒜\textsubscript{X}(X) be the Dolbeault algebra of X and O\textsubscript{X} its sheaf of holomorphic functions. Assume that X satisfies the Dolbeault–Grothendieck lemma (H1) and the local gauge lemma in its weak form (H2). Then the functor J : Ho(P\textsubscript{A}) ⥤ D(O\textsubscript{X}), M ↦ α(M), into the derived category of sheaves of O\textsubscript{X}-modules is full and faithful, and its essential image is exactly D\textsubscript{perf}(X, O\textsubscript{X}). Equivalently, J corestricts to an equivalence Ho(P\textsubscript{A}) ≃ D\textsubscript{perf}(X, O\textsubscript{X}). Both hypotheses are theorems (H1, H2), so this form is a corollary of Z5.j.

Deviation: D\textsubscript{perf}(X, O\textsubscript{X}) instead of Block's D\textsuperscript{b\_}\{coh\}(X); the equivalence is stated as full, faithful and essential image D\textsubscript{perf}.
\par\smallskip{\footnotesize Source: blueprint.md §1.}
\end{theorem}
\begin{proof}
\leanok
\uses{F15}
See the Lean proof and the cited items.
\paragraph{Blueprint notes.}
For a finite-dimensional complex normed space E and a compact Hausdorff complex manifold X modelled on E, let A = 𝒜\textsubscript{X}(X) be the Dolbeault algebra (F10) and O\textsubscript{X} the sheaf of holomorphic functions (F12). Assume the two analytic hypotheses H1 (Dolbeault–Grothendieck) and H2 (local gauge lemma, weak form) on X. Then the functor

J : Ho(P\textsubscript{A}) ⥤ D\textsubscript{perf}(X, O\textsubscript{X}),  M ↦ α(M)

of C1 is an equivalence of categories.

Lean shape (names indicative, fixed in M0):

\begin{verbatim}
theorem blockPerf
    (E : Type) [NormedAddCommGroup E] [NormedSpace ℂ E] [FiniteDimensional ℂ E]
    (X : Type) [TopologicalSpace X] [T2Space X] [CompactSpace X]
    [ChartedSpace E X] [IsManifold 𝓘(ℂ, E) ∞ X]
    (h₁ : DolbeaultGrothendieck E X) (h₂ : LocalGaugeLemma E X) :
    (blockFunctor E X).Full ∧ (blockFunctor E X).Faithful ∧
      (blockFunctor E X).essImage = Perf (ringOf (holSheaf E X))
\end{verbatim}
Here \texttt{blockFunctor E X : HoP (dolbeault E X) ⥤ DerivedCategory (Mod (ringOf (holSheaf E X)))} is J with values in the whole derived category, and \texttt{Perf (ringOf (holSheaf E X))} is the \texttt{ObjectProperty} of S9.c (v1.9: \texttt{Perf} takes the ring sheaf; v1.10: \texttt{Perf} and \texttt{Mod} take a \texttt{RingSheaf X}, while \texttt{holSheaf E X} is a sheaf of \texttt{CommRingCat}, hence \texttt{ringOf}; \texttt{E} is explicit everywhere because the model space of \texttt{ChartedSpace} is not an \texttt{outParam}). \texttt{blockFunctor} and \texttt{Perf} use the same derived-category instance, \texttt{HasDerivedCategory.standard}, fixed once in \texttt{Statements.lean}. \warn{} **Why not \texttt{IsEquivalence} onto D\textsubscript{perf}.** That J lands in D\textsubscript{perf} is C2, which uses LL = H2, so a functor \texttt{Ho(P\_A) ⥤ D\_perf(X)} cannot be defined before H2 is assumed. Fully faithful with essential image exactly D\textsubscript{perf} is equivalent to "the corestriction is an equivalence"; that corollary is proved in \texttt{Checks.lean} (K6).

\warn{} \textbf{Smoothness exponents.} In Mathlib \texttt{⊤ : ℕ∞ω} is ω (analytic) and \texttt{∞} is C\textsuperscript{∞}. Complex manifolds use \texttt{∞} (complex C\textsuperscript{∞} transition maps, i.e. holomorphic); all forms and partitions of unity are real \texttt{∞}. Never write \texttt{⊤} for smoothness in this project.

\textbf{Relation to the sources.} This is [B, Thm 4.1.3] with D\textsuperscript{b\_coh} replaced by D\textsubscript{perf}, and [A, Thm 7.4.3] at h = 0 for compact X. The step D\textsubscript{perf}(X) = D\textsuperscript{b\_coh}(X) is Z1 (LATER).

\textbf{What is deliberately not here} (all in §6): D\textsuperscript{b\_coh} (Z1), triangulated structure (Z2), twisted sheaves and curvature (Z3), proofs of H1 and H2 (Z4, Z5), the dg-level quasi-equivalence (Z6), open manifolds (Z7), comparison with intrinsic forms (Z8).
\end{proof}

\begin{theorem}[F1 — Realification]
\label{F1}
\lean{Cohesive.isManifold_real, Cohesive.contMDiffRing_real_complex}
\leanok
Let E be a complex normed space and M a complex C\textsuperscript{∞} manifold modelled on E (charts with complex C\textsuperscript{∞}, i.e. holomorphic, transition maps). Then M, with the same charts, is a real C\textsuperscript{∞} manifold modelled on E (E regarded as a real normed space). Moreover ℂ, as a real manifold, is a smooth ring: multiplication ℂ × ℂ → ℂ is real C\textsuperscript{∞}. Hence the real-smooth ℂ-valued functions on M form a sheaf of commutative rings C\textsuperscript{∞\_M}(ℂ).
\par\smallskip{\footnotesize Source: blueprint.md §5.}
\end{theorem}
\begin{proof}
\leanok
See the Lean proof and the cited items.
\paragraph{Blueprint notes.}
\texttt{IsManifold 𝓘(ℝ, E) ∞ M} from \texttt{IsManifold 𝓘(ℂ, E) ∞ M}, same charts. Proof: transition maps are \texttt{ContDiffOn ℂ ∞}, hence \texttt{ContDiffOn ℝ ∞} (\texttt{ContDiffOn.restrict\_scalars} ✓). \texttt{NormedSpace ℝ E} comes from the instance \texttt{NormedSpace.complexToReal} ✓. The realification lemma itself is ✗ in Mathlib. Provide it as a theorem plus local instance; avoid a global instance (conventions §7).
\end{proof}

\begin{theorem}[F2 — The fibre Λ\textsuperscript{0,q}]
\label{F2}
\lean{Cohesive.Conj.map, Cohesive.form0Equiv, Cohesive.extToFibre, Cohesive.extToFibre_bijective, Cohesive.extEquivFibre, Cohesive.dzbarBasis}
\leanok
Let E and F be complex normed spaces. The conjugate space Ē is E with the same norm and the scalar action c · v := c̄ v. For q ∈ ℕ, the fibre Λ\textsuperscript{0,q}(E; F) is the normed space of continuous ℂ-multilinear alternating maps Ē\textsuperscript{q} → F (equivalently: real alternating maps on E that are antilinear in each slot).

\begin{itemize}
\item A continuous ℂ-linear map T : E → E′ is also a continuous ℂ-linear map Ē → Ē′.
\item Λ\textsuperscript{0,0}(E; F) ≅ F, isometrically: a 0-form at a point is a value.
\item If E is finite-dimensional, Φ\textsubscript{q} : ⋀\textsuperscript{q\_ℂ} (Ē)\textsuperscript{*} → Λ\textsuperscript{0,q}(E; ℂ), φ\textsubscript{1} ∧ ⋯ ∧ φ\textsubscript{q} ↦ (v ↦ det(φ\textsubscript{j}(v\textsubscript{i}))), is a linear isomorphism.
\item For a basis b of E indexed by a finite linearly ordered set ι, the forms dz̄\textsubscript{I} = Φ\textsubscript{q}(dz̄\textsubscript{i\_1} ∧ ⋯ ∧ dz̄\textsubscript{i\_q}) (dz̄\textsubscript{i} the dual basis of b in Ē), I running over the q-element subsets of ι, form a basis of Λ\textsuperscript{0,q}(E; ℂ). So dim Λ\textsuperscript{0,q}(E; ℂ) = C(n, q), n = dim\textsubscript{ℂ} E, and Λ\textsuperscript{0,q} = 0 for q > n.
\end{itemize}
\par\smallskip{\footnotesize Source: blueprint.md §5.}
\end{theorem}
\begin{proof}
\leanok
See the Lean proof and the cited items.
\paragraph{Blueprint notes.}
Ē = \texttt{Conj E}: E with the conjugate ℂ-action, a type synonym declared with \texttt{def} (not \texttt{abbrev}), normed like E. ✗ in Mathlib. A ℂ-linear map on Ē is an antilinear map on E, and a ℂ-linear map of E (for example DG(x) for holomorphic G) is also ℂ-linear on Ē.

Λ\textsuperscript{0,q}(E; F) := \texttt{Conj E [⋀\textasciicircum{}Fin q]→L[ℂ] F}, a ℂ-normed space (\texttt{instNormedAddCommGroup}, \texttt{instNormedSpace} ✓). \warn{} This replaces the earlier subtype of ℝ-alternating maps antilinear in each slot: the same maps, without membership side conditions, and Mathlib's exterior-power API applies to Ē directly. P2 confirmed \texttt{Conj E} works without instance diamonds; no fallback.

\begin{itemize}
\item \textbf{F2.a} Λ\textsuperscript{0,0}(E; F) ≅ F (\texttt{constOfIsEmpty} ✓).
\item \textbf{F2.b Exterior-algebra bridge.} Φ\textsubscript{q} : ⋀\textsuperscript{q\_ℂ}(Dual\textsubscript{ℂ} Ē) ≃ Λ\textsuperscript{0,q}(E; ℂ), φ₁ ∧ ⋯ ∧ φ\textsubscript{q} ↦ (v ↦ det(φ\textsubscript{i}(v\textsubscript{j}))), from \texttt{exteriorPower.pairingDual} ✓ (\texttt{pairingDual\_ιMulti\_ιMulti} ✓) and \texttt{exteriorPower.alternatingMapLinearEquiv} ✓, plus continuity in finite dimension. Bijective via the basis lemmas \texttt{ιMultiDual\_apply\_diag}, \texttt{ιMultiDual\_apply\_nondiag} ✓ and \texttt{Module.Basis.exteriorPower} ✓; packaging it as an equivalence is ✗.
\item \textbf{F2.c} Λ\textsuperscript{0,q} = 0 for q > n; dim = C(n, q)·dim F; a basis of E gives the basis dz̄\textsubscript{I}.
\end{itemize}
\end{proof}

\begin{theorem}[F3 — The ∂̄-operator on open subsets of E]
\label{F3}
\lean{Cohesive.antiPart, Cohesive.dbarPart, Cohesive.dbar, Cohesive.dbar_add, Cohesive.dbar_const_smul, Cohesive.dbar_congr}
\leanok
Let E, F be complex normed spaces and q ∈ ℕ. The antilinear part of a real-linear map L : E → G is L″(v) = ½(L v + i·L(i v)), a ℂ-linear map Ē → G. For α : E → Λ\textsuperscript{0,q}(E; F) and x ∈ E let D″α(x) be the antilinear part of the real derivative Dα(x), and ∂̄α(x)(v\textsubscript{0}, …, v\textsubscript{q}) = Σ\textsubscript{i} (−1)\textsuperscript{i} D″α(x)(v\textsubscript{i})(v\textsubscript{0}, …, v̂\textsubscript{i}, …, v\textsubscript{q}) ∈ Λ\textsuperscript{0,q+1}(E; F) (Mathlib's normalization of the exterior derivative). Properties:

\begin{itemize}
\item ∂̄(α + β)(x) = ∂̄α(x) + ∂̄β(x) and ∂̄(cα)(x) = c ∂̄α(x) (c ∈ ℂ) if α, β are real-differentiable at x;
\item ∂̄α(x) depends only on the germ of α at x.
\end{itemize}
Deviation: ∂̄ is defined for every function on E (with value 0 where α is not differentiable); by locality this gives ∂̄ on open subsets of E.
\par\smallskip{\footnotesize Source: blueprint.md §5.}
\end{theorem}
\begin{proof}
\leanok
\uses{F5}
See the Lean proof and the cited items.
\paragraph{Blueprint notes.}
For ω : E → Λ\textsuperscript{0,q}(E; F), smooth near x: D″ω(x) = ½ (Dω(x) + i·Dω(x) ∘ J), the antilinear part of the real derivative (J = multiplication by i), regarded as an element of \texttt{Conj E →L[ℂ] Λ\textasciicircum{}\{0,q\}(E; F)}; and ∂̄ω(x) = \texttt{alternatizeUncurryFin} (D″ω(x)) ✓, over ℂ on Ē.

\begin{itemize}
\item \textbf{F3.a} D″ω(x) is ℂ-linear on Ē (the defining computation).
\item \textbf{F3.b} ℂ-linear in ω, local (depends on the germ), smooth in, smooth out.
\item \textbf{F3.c} (sanity) ∂̄ω is the (0, q+1)-part of \texttt{extDeriv} ✓ ω.
\end{itemize}
\end{proof}

\begin{theorem}[F4 — ∂̄∂̄ = 0]
\label{F4}
\lean{Cohesive.dbar_dbar}
\leanok
Let E, F be complex normed spaces and α : E → Λ\textsuperscript{0,q}(E; F). If α is real C² at x, then ∂̄(∂̄α)(x) = 0.
\par\smallskip{\footnotesize Source: blueprint.md §5.}
\end{theorem}
\begin{proof}
\leanok
\uses{F3}
See the Lean proof and the cited items.
\paragraph{Blueprint notes.}
For ω C² near x, ∂̄(∂̄ω)(x) = 0. Proof: copy \texttt{extDerivWithin\_extDerivWithin\_apply} ✓ with D″ in place of D, over ℂ on Ē. The ℂ-bilinear map (u, v) ↦ D″D″ω(x)(u, v) on Ē is symmetric (from \texttt{ContDiffAt.isSymmSndFDerivAt} ✓ over ℝ); conclude with \texttt{alternatizeUncurryFin\_alternatizeUncurryFinCLM\_comp\_of\_symmetric} ✓.
\end{proof}

\begin{theorem}[F5 — Product on the total fibre]
\label{F5}
\lean{Cohesive.TotFibre.ofDeg, Cohesive.TotFibre.deg, Cohesive.TotFibre.mulCLM, Cohesive.TotFibre.extEquivTotFibre_mul, Cohesive.TotFibre.ofDeg_mul_ofDeg, Cohesive.TotFibre.ofDeg_mul_comm}
\leanok
Let E be a finite-dimensional complex normed space, n = dim\textsubscript{ℂ} E. The total fibre is the normed space Λ\textsuperscript{0,•}(E) = ∏\textsubscript{q=0}\textsuperscript{n} Λ\textsuperscript{0,q}(E; ℂ). It is a ℂ-algebra, with the product transported from the exterior algebra ⋀\textsubscript{ℂ} (Ē)\textsuperscript{*} along Φ = (Φ\textsubscript{q})\textsubscript{q} of F2, so that Φ is an isomorphism of ℂ-algebras. Viewing a q-form as an element of Λ\textsuperscript{0,•}(E) in degree q (zero if q > n):

\begin{itemize}
\item the product of a p-form and a q-form is a (p+q)-form;
\item graded commutativity: α ∧ β = (−1)\textsuperscript{pq} β ∧ α for a p-form α and a q-form β;
\item Laplace formula: for a 1-form φ and a q-form β, (φ ∧ β)(v\textsubscript{0}, …, v\textsubscript{q}) = Σ\textsubscript{i} (−1)\textsuperscript{i} φ(v\textsubscript{i}) β(v\textsubscript{0}, …, v̂\textsubscript{i}, …, v\textsubscript{q});
\item the product is a continuous bilinear map.
\end{itemize}
\par\smallskip{\footnotesize Source: blueprint.md §5.}
\end{theorem}
\begin{proof}
\leanok
\uses{F2}
See the Lean proof and the cited items.
\paragraph{Blueprint notes.}
The total fibre is Λ\textsuperscript{0,•}(E) := Π\textsubscript{q ≤ n} Λ\textsuperscript{0,q}(E; ℂ), a finite product of normed spaces. Its multiplication is transported along Φ = (Φ\textsubscript{q})\textsubscript{q} from \texttt{ExteriorAlgebra ℂ (Dual\_ℂ Ē)}, which is graded (\texttt{ExteriorAlgebra.gradedAlgebra} ✓) and graded-commutative on generators (\texttt{ιMulti\_mul\_ιMulti\_anticomm} ✓). Associativity, unit and graded commutativity come for free; continuity by finite dimension. \warn{} This replaces a hand-built shuffle product, whose associativity is not in Mathlib (only \texttt{AlternatingMap.domCoprod} ✓ and its alternatization lemmas are).

\begin{itemize}
\item \textbf{F5.a Laplace formula.} For a 1-form φ and a q-form ω, φ ∧ ω = \texttt{alternatizeUncurryFin} (v ↦ φ(v) • ω), i.e. (φ ∧ ω)(v₀, …, v\textsubscript{q}) = Σ\textsubscript{i} (−1)\textsuperscript{i} φ(v\textsubscript{i}) ω(…, v̂\textsubscript{i}, …). From expansion of the determinant along the first row (\texttt{Matrix.det\_succ\_row\_zero} ✓). This is the only link between the product and the analysis.
\item \textbf{F5.b} F-valued forms are a module over Λ\textsuperscript{0,•}(E) via Λ\textsuperscript{0,q}(E; F) ≅ Λ\textsuperscript{0,q}(E; ℂ) ⊗ F (later only).
\end{itemize}
\end{proof}

\begin{theorem}[F6 — Leibniz rule]
\label{F6}
\lean{Cohesive.wedge, Cohesive.dbar_wedge}
\leanok
\uses{F2, F5}
Let E be a finite-dimensional complex normed space. For a p-form α and a q-form β at a point, α ∧ β ∈ Λ\textsuperscript{0,p+q}(E; ℂ) is the degree-(p+q) part of their product in Λ\textsuperscript{0,•}(E). If α : E → Λ\textsuperscript{0,p}(E; ℂ) and β : E → Λ\textsuperscript{0,q}(E; ℂ) are real-differentiable at x, then in Λ\textsuperscript{0,•}(E) ∂̄(α ∧ β)(x) = ∂̄α(x) ∧ β(x) + (−1)\textsuperscript{p} α(x) ∧ ∂̄β(x).
\par\smallskip{\footnotesize Source: blueprint.md §5.}
\end{theorem}
\begin{proof}
\leanok
\uses{F3}
See the Lean proof and the cited items.
\paragraph{Blueprint notes.}
∂̄(α ∧ β) = ∂̄α ∧ β + (−1)\textsuperscript{p} α ∧ ∂̄β. Proof: (i) ∂̄(f·η) = ∂̄f ∧ η + f·∂̄η for a function f, by the product rule and F5.a; (ii) constant forms are ∂̄-closed; (iii) locally every form is Σ\textsubscript{I} f\textsubscript{I} dz̄\textsubscript{I} (F2.c); expand and use graded commutativity. No general alternation identity is needed.
\end{proof}

\begin{theorem}[F7 — Holomorphic pullback]
\label{F7}
\lean{Cohesive.TotFibre.pull, Cohesive.pullback, Cohesive.dbar_pullback, Cohesive.pullback_smul, Cohesive.TotFibre.pull_mul, Cohesive.TotFibre.pull_one}
\leanok
Let E, F be complex normed spaces. For G : E → E and α : E → Λ\textsuperscript{0,q}(E; F) the pullback is G\textsuperscript{*α}(x) = α(G x) ∘ DG(x), with the complex derivative DG(x) acting on Ē (junk where G is not complex-differentiable).

\begin{itemize}
\item If α is real-differentiable at G x and G is complex C² at x, then ∂̄(G\textsuperscript{\emph{α})(x) = G\textsuperscript{}}(∂̄α)(x).
\item G\textsuperscript{*}(f α)(x) = f(G x) · G\textsuperscript{*α}(x) for f : E → ℂ.
\item (G ∘ H)\textsuperscript{\emph{} α (x) = H\textsuperscript{}}(G\textsuperscript{*α})(x) if H is complex-differentiable at x and G at H x.
\item If E is finite-dimensional, the pullback T\textsuperscript{*} along a continuous ℂ-linear T : E → E, applied degreewise on Λ\textsuperscript{0,•}(E), is a unital ring homomorphism.
\end{itemize}
Deviation: G is a map E → E rather than a holomorphic map between open subsets; since all hypotheses are local at x, this covers the latter.
\par\smallskip{\footnotesize Source: blueprint.md §5.}
\end{theorem}
\begin{proof}
\leanok
\uses{F3}
See the Lean proof and the cited items.
\paragraph{Blueprint notes.}
For G : V → V′ holomorphic between open subsets of E, G\textsuperscript{*ω}(x) = ω(G x) ∘ DG(x) (\texttt{compContinuousLinearMap} ✓), with DG(x) regarded as a ℂ-linear map of Ē.

\begin{itemize}
\item \textbf{F7.a} G\textsuperscript{\emph{} ∂̄ = ∂̄ G\textsuperscript{}}. Proof: copy \texttt{extDerivWithin\_pullback} ✓ with D″. The second-derivative term is antilinear in the new slot, so D″ keeps it, and alternation kills it (\texttt{alternatizeUncurryFin\_fderivCompContinuousLinearMap\_eq\_zero} ✓).
\item \textbf{F7.b} On fibres G\textsuperscript{*} is the algebra map induced by the dual of DG(x) on the exterior algebra (\texttt{ExteriorAlgebra.map}), so it respects products; it respects multiplication by functions.
\item \textbf{F7.c} (G ∘ H)\textsuperscript{\emph{} = H\textsuperscript{}} G\textsuperscript{\emph{}, id\textsuperscript{}} = id.
\end{itemize}
\end{proof}

\begin{theorem}[F8 — Cauchy–Riemann]
\label{F8}
\lean{Cohesive.dbar_eq_zero_iff_contDiffOn}
\leanok
Let E, F be complex normed spaces, V ⊆ E open and f : E → F real C\textsuperscript{∞} on V. Then ∂̄f = 0 at every point of V (f regarded as a 0-form) if and only if f is complex C\textsuperscript{∞} (holomorphic) on V.
\par\smallskip{\footnotesize Source: blueprint.md §5.}
\end{theorem}
\begin{proof}
\leanok
\uses{F3}
See the Lean proof and the cited items.
\paragraph{Blueprint notes.}
For f : E → F real C\textsuperscript{∞} on an open V: ∂̄f = 0 on V ⇔ f is \texttt{ContDiffOn ℂ ∞} on V. Proof (⇒): ∂̄f = 0 means Df(x) is ℂ-linear (\texttt{differentiableAt\_iff\_restrictScalars} ✓). A symmetric ℝ-bilinear map that is ℂ-linear in one slot is ℂ-linear in both, so the derivative of fderiv ℂ f is again ℂ-linear; induct. \warn{} This avoids several-variable Cauchy or Osgood theory, which Mathlib does not have (✗ "Osgood", "Hartogs").
\end{proof}

\begin{theorem}[F9 — The sheaves 𝒜\textsuperscript{0,q}\textsubscript{M}(F) and 𝒜\textsuperscript{0,•}\textsubscript{M}]
\label{F9}
\lean{Cohesive.IsDolbeault, Cohesive.IsOfDegree.dbarSec, Cohesive.restrictForm, Cohesive.dolbeaultLocalPredicate, Cohesive.IsDolbeault.contDiffAt_chartExpr}
\leanok
\uses{F7}
Let E be a finite-dimensional complex normed space and X a complex manifold modelled on E. A form on an open U ⊆ X is a function α : U → Λ\textsuperscript{0,•}(E), the value α(y) being written in the chart at y. Its expression in a chart e is z ↦ T\textsuperscript{\emph{α}(y), y = e\textsuperscript{−1}(z), T the complex derivative at z of the transition map chartAt(y) ∘ e\textsuperscript{−1}. α is a *Dolbeault form} if, for every x ∈ U, its expression in φ = chartAt(x) is real C\textsuperscript{∞} at φ(x); its ∂̄ is (∂̄α)(x) = ∂̄(expression of α in φ)(φ(x)), degree by degree. A form has degree q if all its other components vanish.

\begin{itemize}
\item Restrictions of Dolbeault forms are Dolbeault forms, and being Dolbeault is local, so Dolbeault forms are a sheaf of sets on X.
\item The expression of a Dolbeault form in any atlas chart e is real C\textsuperscript{∞} at e(x) for x ∈ U ∩ e.source.
\item ∂̄ of a Dolbeault form is a Dolbeault form, and ∂̄ maps forms of degree q to degree q + 1.
\end{itemize}
Deviation: forms take values in the total fibre, and only ℂ-valued forms are treated; 𝒜\textsuperscript{0,q} is the degree-q part (equivalent, since the direct sum is finite).
\par\smallskip{\footnotesize Source: blueprint.md §5.}
\end{theorem}
\begin{proof}
\leanok
\uses{F4, F6}
See the Lean proof and the cited items.
\paragraph{Blueprint notes.}
Sections over an open U are functions ω : U → Λ\textsuperscript{0,q}(E; F) (or into the total fibre Λ\textsuperscript{0,•}(E)), ω(x) read in the chart \texttt{chartAt x}, satisfying the local predicate: for each x ∈ U, with φ = \texttt{chartAt x}, the chart expression z ↦ ω(y) ∘ D(\texttt{chartAt y} ∘ φ⁻¹)(z) at z = φ(y) is smooth near φ(x). Built with \texttt{TopCat.LocalPredicate} ✓ and \texttt{TopCat.subsheafToTypes} ✓, as Mathlib builds \texttt{smoothSheaf} ✓; the ring structure on 𝒜\textsuperscript{0,•} follows \texttt{smoothSheafCommRing} ✓. Making each 𝒜\textsuperscript{0,q} a \texttt{SheafOfModules} over 𝒜\textsuperscript{0} has no Mathlib precedent and is part of F9's size.

\begin{itemize}
\item \textbf{F9.a} The predicate is independent of the chart choices (F7). The transition factor is the coordinate change of \texttt{tangentBundleCore} ✓ (\texttt{tangentBundleCore\_coordChange} ✓), so its cocycle lemmas apply.
\item \textbf{F9.b} ∂̄ chartwise: (∂̄ω)(x) = (∂̄ω\textsubscript{φ})(φ(x)), φ = \texttt{chartAt x}; well defined by F7.a.
\item \textbf{F9.c} ∧ and the module structure over smooth functions are pointwise.
\item \textbf{F9.d} 𝒜\textsuperscript{0,0}\textsubscript{M}(ℂ) ≅ C\textsuperscript{∞\_M}(ℂ) := \texttt{smoothSheafCommRing 𝓘(ℝ, E) 𝓘(ℝ, ℂ) M ℂ} ✓.
\item \textbf{F9.e} (sanity) For an open V ⊆ E with its induced structure (\texttt{TopologicalSpace.Opens.instChartedSpace} ✓), sections over V are just smooth functions V → Λ\textsuperscript{0,q}(E; F).
\end{itemize}
\end{proof}

\begin{theorem}[F10 — The Dolbeault dga]
\label{F10}
\lean{Cohesive.dolbeault0, Cohesive.dolbeaultSubalgebra, Cohesive.dbarL, Cohesive.dolbeaultRes, Cohesive.dolbeault_glue, Cohesive.isOfDegree_of_locally}
\leanok
\uses{A1, F5, F9}
Let E be a finite-dimensional complex normed space, n = dim\textsubscript{ℂ} E, and X a complex manifold modelled on E. For an open U ⊆ X, 𝒜\textsubscript{X}(U) is the ℂ-algebra of Dolbeault forms on U (F9) with pointwise operations, graded by 𝒜\textsuperscript{0,q}\textsubscript{X}(U) = the forms of degree q (zero for q < 0 and q > n).

\begin{itemize}
\item 𝒜\textsubscript{X}(U) with d = ∂̄ is a bounded graded-commutative dga (A1) with top degree n and curvature 0.
\item Restriction along V ⊆ U is a map of dgas 𝒜\textsubscript{X}(U) → 𝒜\textsubscript{X}(V).
\item Sheaf conditions: Dolbeault forms on opens U\textsubscript{i} that agree on overlaps glue to a unique Dolbeault form on ⋃ U\textsubscript{i}; having degree q is a local property.
\item U ↦ 𝒜\textsuperscript{0,0}\textsubscript{X}(U) is a sheaf 𝒜\textsuperscript{0} of commutative rings on X.
\end{itemize}
The Dolbeault algebra of X is A = 𝒜\textsubscript{X}(X).
\par\smallskip{\footnotesize Source: blueprint.md §5.}
\end{theorem}
\begin{proof}
\leanok
\uses{A9}
See the Lean proof and the cited items.
\paragraph{Blueprint notes.}
𝒜\textsubscript{M}(U) := sections over U of the total-fibre sheaf 𝒜\textsuperscript{0,•}\textsubscript{M} (F9), with pointwise product (F5) and ∂̄ (F9.b); internally graded by the components q, and 𝒜\textsuperscript{0,q}\textsubscript{M} is the q-th component sheaf. A bounded graded-commutative dga (A1) for every open U, with restrictions dga maps. (v1.9) The component sheaves are \emph{defined} as the subtypes 𝒜\textsuperscript{0,q}\textsubscript{M}(U) := ↥(𝒜\textsubscript{M}(U) q) of the total section algebra, so that the 𝒜\textsuperscript{0} of C0 is the dga's own \texttt{𝒜 0} (C0, one carrier for A⁰). \textbf{The Dolbeault algebra of X is A = 𝒜\textsubscript{X}(X).} \warn{} The total fibre removes the direct-sum construction of conventions §5: the ring axioms hold pointwise.
\end{proof}

\begin{theorem}[F11 — Fineness and local freeness]
\label{F11}
\lean{Cohesive.isFine_dolbeault0, Cohesive.dolbeault_locallyFree}
\leanok
\uses{F2}
Let E be a finite-dimensional complex normed space, n = dim\textsubscript{ℂ} E, and X a complex manifold modelled on E.

\begin{itemize}
\item If X is compact Hausdorff, 𝒜\textsuperscript{0} is fine: for every finite open cover (U\textsubscript{i}) of X there are global sections ρ\textsubscript{i} ∈ 𝒜\textsuperscript{0,0}\textsubscript{X}(X) with Σ\textsubscript{i} ρ\textsubscript{i} = 1, each ρ\textsubscript{i} supported in U\textsubscript{i} (it vanishes on an open V\textsubscript{i} with U\textsubscript{i} ∪ V\textsubscript{i} = X).
\item 𝒜\textsuperscript{0,q} is locally free of rank C(n, q) over 𝒜\textsuperscript{0}: every x ∈ X has an open neighbourhood U with forms e\textsubscript{I} ∈ 𝒜\textsuperscript{0,q}\textsubscript{X}(U), I running over the q-element subsets of \{0, …, n−1\}, such that on every open V ⊆ U each β ∈ 𝒜\textsuperscript{0,q}\textsubscript{X}(V) is uniquely Σ\textsubscript{I} f\textsubscript{I} · e\textsubscript{I}|\textsubscript{V} with f\textsubscript{I} ∈ 𝒜\textsuperscript{0,0}\textsubscript{X}(V).
\end{itemize}
\par\smallskip{\footnotesize Source: blueprint.md §5.}
\end{theorem}
\begin{proof}
\leanok
\uses{F1, F12, S1}
See the Lean proof and the cited items.
\paragraph{Blueprint notes.}
For compact X: 𝒜\textsuperscript{0} is fine (\texttt{SmoothPartitionOfUnity.exists\_isSubordinate} ✓ for 𝓘(ℝ, E); needs T2 and σ-compact); 𝒜\textsuperscript{0,q} is locally free of rank C(n, q) over 𝒜\textsuperscript{0} (F2.c in charts); 𝒜\textsuperscript{0,q} = 0 for q > n. (v1.9) \texttt{IsFine} asks for the stalk-free \texttt{SupportedIn} (∃ V, U ⊔ V = ⊤ ∧ ρ|V = 0): take V := (tsupport ρ\textsubscript{i})ᶜ; the real-valued ρ\textsubscript{i} become sections of 𝒜\textsuperscript{0} through ℝ → ℂ.
\end{proof}

\begin{theorem}[F12 — Holomorphic functions]
\label{F12}
\lean{Cohesive.dolbeaultRhoU, Cohesive.dolbeaultRho, Cohesive.dolbeaultRhoU_injective, Cohesive.dbar_dolbeaultRhoU, Cohesive.exists_dolbeaultRhoU_of_dbar_eq_zero, Cohesive.holSheaf}
\leanok
\uses{F8}
Let E be a finite-dimensional complex normed space and X a complex manifold modelled on E. O\textsubscript{X} is the sheaf of commutative rings of complex C\textsuperscript{∞} (holomorphic) functions on X. ρ : O\textsubscript{X} → 𝒜\textsuperscript{0} is the morphism of sheaves of rings that sends a holomorphic function f on U to f as a form of degree 0. For every open U ⊆ X:

\begin{itemize}
\item ρ\textsubscript{U} is injective;
\item ∂̄ ρ\textsubscript{U}(f) = 0 for every f ∈ O\textsubscript{X}(U);
\item every α ∈ 𝒜\textsuperscript{0,0}\textsubscript{X}(U) with ∂̄α = 0 is ρ\textsubscript{U}(f) for some f ∈ O\textsubscript{X}(U).
\end{itemize}
So 0 → O\textsubscript{X}(U) → 𝒜\textsuperscript{0,0}\textsubscript{X}(U) → 𝒜\textsuperscript{0,1}\textsubscript{X}(U) is exact.
\par\smallskip{\footnotesize Source: blueprint.md §5.}
\end{theorem}
\begin{proof}
\leanok
\uses{F10}
See the Lean proof and the cited items.
\paragraph{Blueprint notes.}
O\textsubscript{M} := \texttt{smoothSheafCommRing 𝓘(ℂ, E) 𝓘(ℂ, ℂ) M ℂ} ✓ (complex C\textsuperscript{∞} = holomorphic). ρ : O\textsubscript{M} → C\textsuperscript{∞\_M}(ℂ) is the inclusion (F1 and restriction of scalars); it is injective with image ker(∂̄ : 𝒜\textsuperscript{0} → 𝒜\textsuperscript{1}) by F8. So 0 → O → 𝒜\textsuperscript{0} → 𝒜\textsuperscript{1} is exact.
\end{proof}

\begin{theorem}[H1 — Dolbeault–Grothendieck lemma]
\label{H1}
\lean{Cohesive.DolbeaultGrothendieck, Cohesive.dolbeaultGrothendieck}
\leanok
Let E be a finite-dimensional complex normed space and X a complex manifold modelled on E. X satisfies the \emph{Dolbeault–Grothendieck lemma} if for every open U ⊆ X, every q ∈ ℕ, every x ∈ U and every α ∈ 𝒜\textsuperscript{0,q+1}\textsubscript{X}(U) with ∂̄α = 0, there are an open U′ with x ∈ U′ ⊆ U and η ∈ 𝒜\textsuperscript{0,q}\textsubscript{X}(U′) with ∂̄η = α|\textsubscript{U′}. Theorem: this holds for every such X (no compactness or Hausdorff assumption). With F12 it says that ρ : O\textsubscript{X} → 𝒜\textsuperscript{0,•}\textsubscript{X} is a resolution.

Deviation: stated on the manifold X in sheaf form, not on polydiscs in ℂ\textsuperscript{n}.
\par\smallskip{\footnotesize Source: blueprint.md §5.}
\end{theorem}
\begin{proof}
\leanok
\uses{F10}
See the Lean proof and the cited items.
\paragraph{Blueprint notes.}
\texttt{DolbeaultGrothendieck E X : Prop}: for every open U ⊆ X, q ∈ ℕ, x ∈ U and ω ∈ Γ(U, 𝒜\textsuperscript{0,q+1}\textsubscript{X}) with ∂̄ω = 0, there are an open x ∈ U′ ⊆ U and η ∈ Γ(U′, 𝒜\textsuperscript{0,q}\textsubscript{X}) with ∂̄η = ω|U′. Together with F12 this says exactly that ρ : O\textsubscript{X} → 𝒜\textsuperscript{•\_X} is a resolution.

\begin{itemize}
\item True on every complex manifold (Dolbeault–Grothendieck lemma). Proof is Z4.
\item Trivial cases to prove as lemmas: q + 1 > n; n = 0.
\item \warn{} Stated on X, not on open subsets of E: this keeps chart transport (F13) off the critical path, and it is the standard sheaf-theoretic form of the lemma.
\end{itemize}
\end{proof}

\begin{theorem}[H2 — Local gauge lemma, weak form]
\label{H2}
\lean{Cohesive.LocalGaugeLemma, Cohesive.localGaugeLemma}
\leanok
\uses{A10, A9}
Let E be a finite-dimensional complex normed space and X a complex manifold modelled on E. X satisfies the \emph{local gauge lemma} (weak form) if for every cohesive module M over A = 𝒜\textsubscript{X}(X) and every x ∈ X there are an open U′ ∋ x, a bounded complex W of finitely generated projective O\textsubscript{X}(U′)-modules, and an isomorphism in Ho(P\textsubscript{𝒜\_X(U′)}) between the restriction M|\textsubscript{U′} (base change along 𝒜\textsubscript{X}(X) → 𝒜\textsubscript{X}(U′), A9) and 𝒜\textsubscript{X}(U′) ⊗\textsubscript{O\_X(U′)} W (A10, along ρ\textsubscript{U′}). Theorem: this holds for every such X.

Deviation: the weak, sectionwise form of [B, Lemma 4.1.5]: only restrictions of global cohesive modules, a homotopy equivalence instead of a gauge equivalence, and W need not be free.
\par\smallskip{\footnotesize Source: blueprint.md §5.}
\end{theorem}
\begin{proof}
\leanok
\uses{F12}
See the Lean proof and the cited items.
\paragraph{Blueprint notes.}
\texttt{LocalGaugeLemma E X : Prop}: for every cohesive module M over A = 𝒜\textsubscript{X}(X) and x ∈ X, there are an open x ∈ U′, a bounded complex W of finitely generated projective O\textsubscript{X}(U′)-modules, and a homotopy equivalence in P\textsubscript{𝒜\_X(U′)} between M|U′ (A9) and 𝒜\textsubscript{X}(U′) ⊗\textsubscript{O\_X(U′)} W (A10).

\begin{itemize}
\item True: the proof of [B, Lemma 4.1.5] (whose statement is full faithfulness of α) gives a gauge equivalence on a smaller polydisc with W free; this is its weak form, and [A, Assumption 6] at the level of sections. It is field LL of C0 verbatim. Proof is Z5.
\item Trivial case to prove: n = 0.
\item (v1.9, audit F9) Lean form, to be reviewed at gate 3: \texttt{LocalGaugeLemma E X} says: for every \texttt{M : CohesiveModule (dolbeault E X)} and x, there are an open U′ ∋ x, a complex \texttt{W : CochainComplex (ModuleCat.\{0\} (O(U′))) ℤ} with \texttt{IsBoundedFGProj}, and an isomorphism in \texttt{HoP} of the U′-dga between \texttt{M.baseChange (res\_\{U′\})} (A9) and \texttt{CohesiveModule.fromComplex ρ\_\{U′\} … W} (A10, with \texttt{hc : c = 0} and \texttt{d ∘ ρ\_\{U′\} = 0}, both facts about the Dolbeault dga). ρ\textsubscript{U′} : O(U′) →+* ↥(𝒜\textsubscript{U′} 0) (C0, one carrier).
\end{itemize}
\end{proof}

\begin{definition}[F14 — The fields of C0]
\label{F14}
\lean{Cohesive.dolbeaultPresheaf, Cohesive.dolbeaultPreData, Cohesive.dolbeaultData}
\leanok
\uses{C0, F10, F11, F12, H1, H2}
Let E be a finite-dimensional complex normed space and X a complex manifold modelled on E.

\begin{itemize}
\item U ↦ 𝒜\textsubscript{X}(U) with restriction of forms is a presheaf of dgas on X.
\item If X is compact Hausdorff, the Dolbeault data form pre-resolving data (C0.a): the dgas 𝒜\textsubscript{X}(U), which glue and have curvature 0; O = O\textsubscript{X} with ρ\textsubscript{U} : O\textsubscript{X}(U) → 𝒜\textsuperscript{0,0}\textsubscript{X}(U) and ∂̄ ∘ ρ = 0 (F12); 𝒜\textsuperscript{0} a fine sheaf and each 𝒜\textsuperscript{0,q} locally free over it (F11). No analytic hypothesis is used.
\item If moreover X satisfies H1 and H2, they form resolving data (C0.b): ρ is injective, O\textsubscript{X} → 𝒜\textsuperscript{0} → 𝒜\textsuperscript{1} is exact (F12), the Dolbeault complex is locally exact in degrees ≥ 1 (H1), and the local lifting field LL is H2.
\end{itemize}
\par\smallskip{\footnotesize Source: blueprint.md §5.}
\end{definition}
\paragraph{Blueprint notes.}
F10, F11 and F12 give the pre-resolving data C0.a for X, unconditionally. Given H1 and H2, F12 and H1 give the resolution field and H2 is the field LL, completing C0.b.

\begin{theorem}[F15 — Main theorem conditional: proof]
\label{F15}
\lean{Cohesive.blockFunctor, Cohesive.blockPerf}
\leanok
\uses{C0, C6, F14, H1, H2}
The proof of M, as the Lean theorem \texttt{blockPerf}. Let E be a finite-dimensional complex normed space and X a compact Hausdorff complex manifold modelled on E, A = 𝒜\textsubscript{X}(X). The Dolbeault data of X are pre-resolving data (F14) without any analytic hypothesis, so C1 gives the functor J : Ho(P\textsubscript{A}) ⥤ D(O\textsubscript{X}). If X satisfies H1 and H2, the Dolbeault data are resolving data (F14), and the comparison theorem C6 specialized to them says that J is full and faithful with essential image D\textsubscript{perf}(X, O\textsubscript{X}).

This is the only node that applies the abstract theorem C6 to the Dolbeault data; M is the statement itself (\texttt{MainStatementConditional}), and Z5.j is its form without H1 and H2.

Deviation: proved for any model of the derived category (any \texttt{HasDerivedCategory} instance), not only the standard one.
\par\smallskip{\footnotesize Source: blueprint.md §5.}
\end{theorem}
\begin{proof}
\leanok
See the Lean proof and the cited items.
\paragraph{Blueprint notes.}
\texttt{blockFunctor E X := J} for the C0.a data of F14 (no hypotheses). Assemble C0.b from H1, H2 (F14) and apply C6. This proves M.
\end{proof}
